% !TeX root = ../QuestionSet.tex
% ==================== 第8章 导数及其应用 ====================
\QuestionSetChapter{导数及其应用}

% ===== 8.1 导数的概念与运算 =====
\QuestionSetSection{导数的概念与运算}

\questionGroup{一、选择题}

% ----------------------------------------
\begin{question}
若直线$y=2x+5$是曲线$y=e^x+x+a$的切线，则$a=$\blank{3cm}．
\end{question}
\begin{answer}{4}
由 $y'=e^x+1=2$ 得 $x=0$，代入切线 $y=2x+5$ 得 $y=5$，再代入曲线 $y=e^x+x+a$ 得 $5=1+0+a\Rightarrow a=4$。
\end{answer}

% ----------------------------------------
\begin{question}
（1）设函数$f(x)=5\cos x-\cos 5x$,求$f(x)$在$\left[ 0,\dfrac{\pi}{4} \right]$的最大值；

（2）给定$\theta \in (0,\pi)$，设a为实数，证明：存在$y\in [a-\theta ,a+\theta ]$，使得$\cos y\le \cos \theta $；

（3）若存在$\varphi $使得对任意x，都有$5\cos x-\cos (5x+\varphi )\le b$，求b的最小值．
\end{question}
\begin{answer}{(1) $3\sqrt3$；(2) 略；(3) $3\sqrt3$。}
(1) $f'(x)=5(\sin5x-\sin x)=10\cos3x\sin2x$，令 $f'(x)=0$ 得 $x=\dfrac{\pi}{6}$，$f(x)_{\max}=3\sqrt3$。(2) 令 $g(y)=\cos y-\cos\theta$，则 $g(a-\theta)g(a+\theta)\le0$，故存在 $y$ 使 $\cos y\le\cos\theta$。(3) $h(x)=5\cos x-\cos(5x+\varphi)$，分析知 $h(x)_{\max}=3\sqrt3$，故 $b_{\min}=3\sqrt3$。
\end{answer}

\begin{question}[GPT生成]
函数$f(x)=x^3-3x$的导函数为\choiceBox
\choices{$3x^2-3$}{$3x^2+3$}{$x^2-3$}{$3x-3$}
\end{question}
\begin{answer}{A}
由求导公式得 $(x^3)'=3x^2$，$(-3x)'=-3$，故 $f'(x)=3x^2-3$。选 A。
\end{answer}

\questionGroup{二、填空题}

% ----------------------------------------
\begin{question}[GPT生成]
函数$f(x)=\ln x$在$x=1$处的切线斜率为\blank{3cm}．
\end{question}
\begin{answer}{1}
$f'(x)=\dfrac1x$，所以 $f'(1)=1$，即切线斜率为 $1$。
\end{answer}

% ===== 8.2 导数与函数的单调性、极值 =====
\QuestionSetSection{导数与函数的单调性、极值}

\questionGroup{三、解答题}

% ----------------------------------------
\begin{question}[GPT生成]
已知$f(x)=x^3-3x^2+1$，求函数$f(x)$的单调区间。
\end{question}
\begin{answer}{增区间 $(-\infty,0)$、$(2,+\infty)$；减区间 $(0,2)$}
$f'(x)=3x^2-6x=3x(x-2)$。当 $x<0$ 或 $x>2$ 时，$f'(x)>0$，函数单调递增；当 $0<x<2$ 时，$f'(x)<0$，函数单调递减。因此增区间为 $(-\infty,0)$、$(2,+\infty)$，减区间为 $(0,2)$。
\end{answer}

% ----------------------------------------
\begin{question}[GPT生成]
函数$f(x)=x^3-3x$在区间$[-2,2]$上的最大值为\blank{3cm}．
\end{question}
\begin{answer}{2}
$f'(x)=3x^2-3=3(x-1)(x+1)$，临界点为 $x=\pm1$。计算 $f(-2)=-2$，$f(-1)=2$，$f(1)=-2$，$f(2)=2$，故最大值为 $2$。
\end{answer}

\QuestionSetSection*{本章综合训练}

\questionGroup{综合解答题}

% ----------------------------------------
\begin{question}[GPT生成]
已知函数$f(x)=x^3-3x^2+1$。（1）求$f(x)$的单调区间；（2）求$f(x)$的极大值与极小值；（3）求$f(x)$在区间$[-1,3]$上的最大值和最小值。
\end{question}
\solutionSpace{5cm}
\begin{answer}{(1) 增区间$(-\infty,0),(2,+\infty)$，减区间$(0,2)$；(2) 极大值$1$，极小值$-3$；(3) 最大值$1$，最小值$-3$}
$f'(x)=3x^2-6x=3x(x-2)$。当 $x<0$ 或 $x>2$ 时，$f'(x)>0$；当 $0<x<2$ 时，$f'(x)<0$，故单调区间如上。由导数符号变化知 $x=0$ 处取极大值 $f(0)=1$，$x=2$ 处取极小值 $f(2)=-3$。比较 $f(-1)=-3,f(0)=1,f(2)=-3,f(3)=1$，得最大值 $1$，最小值 $-3$。
\end{answer}

% ----------------------------------------
\begin{question}[GPT生成]
设函数$f(x)=e^x+ax$。（1）若曲线在$x=0$处的切线与直线$y=3x+2$平行，求$a$；（2）写出该切线方程；（3）在此$a$下，求不等式$f'(x)>0$的解集。
\end{question}
\solutionSpace{5cm}
\begin{answer}{(1) $a=2$；(2) $y=3x+1$；(3) $\mathbf{R}$}
$f'(x)=e^x+a$。切线斜率为 $f'(0)=1+a$，与 $y=3x+2$ 平行，所以 $1+a=3$，得 $a=2$。此时 $f(0)=1$，切线方程为 $y-1=3(x-0)$，即 $y=3x+1$。又 $f'(x)=e^x+2>0$ 对任意实数 $x$ 恒成立，故解集为 $\mathbf{R}$。
\end{answer}

\printChapterAnswers
